\scanpage{9}
\normalmarginpar
\margin{\handwritten{\unsure{11/4/67|17/4/67|11.4.67|17.4.67}}}
The proof of 7.1. is not correct, when you conte\inserted{n}d on line -9: there remains to be proven the following: ... Already when $D=0$, the proof here would have to introduce conne\unsure{\deleted{f}c|\deleted{c}f}ted étale coverings which are \uline{not Galois} ! This very strongly suggests that a notion of tame ramification should be introduced also \deleted{in} for non Galois coverings. The same remark applies to the proof of 10.1. Maybe you could get along some way in 7.1. using the normali\deleted{\unsure{x|s}}ty assumptions, but I am convienced that these assumptions are anyhow artificial, as well as the assumption that the $D_i$ should be reduced somewhere. You do not seem to make any use of these facts, really
% The right paper edge follows the fully visible y of really; no terminal point is visible.

Also, one feels that 7.\unsure{\deleted{4}5|\deleted{5}4}. should be generalized to the case of a tamely ramified covering, and that it should come out trivially once the generalities have been dealt with properly.

I hope that the theory will come out more clearly and correctly by devoting some care to generalities on ramification data (not necessarily of Kummer type). I will try to write something up within the next days. Please excuse me for the trouble I caused you by not learning my lesson well enough before I put you to work !

It would be very nice indeed to have an appendix on Lefschetz theorem for the tame fundamental group, and it should not be hard to write it. However, if would be safer to wait till the general theory of tame ramification is written up !

Sincerely your's\\
% archive sheet number: 8
